\begin{abstract}
Recent empirical studies on Vision-Language Model (VLM) efficiency argue that autoregressive text decoding dominates 86\%--97\% of total inference energy, leading to the contentious claim that ``seeing is free, speaking is not'' and that visual token reduction yields negligible power savings ($\le$10\%). In this work, we demonstrate that this conclusion is an artifact of task and prompt selection bias---specifically profiling open-ended image captioning where generated response lengths are artificially inflated. In real-world multimodal workloads such as document understanding, chart reasoning, and fine-grained visual probing, the operational profile is completely inverted: extracting key factual entities requires only a few generated tokens ($N_{\text{out}} \in [1, 5]$), whereas resolving small typography and complex visual layout demands dense visual token representations ($N_{\text{vis}} \gg 1000$). We introduce \textbf{GreenVLM}, a task-conditioned thermodynamic framework and benchmark spanning four multimodal task archetypes (DocVQA, ChartQA, POPE, and COCO) across multiple input resolutions ($224^2$, $448^2$, $896^2$). Across 300 systematic evaluations, we show that the inference energy bottleneck completely inverts: in document intelligence and structured probing, visual prefill consumes \textbf{98.9\% to 100.0\%} of total inference energy, rendering decoding cost negligible ($<$1 Joule). Crucially, we construct the first joint Energy--Accuracy Pareto frontier for VLMs, demonstrating that scaling resolution buys up to a \textbf{+55\% relative accuracy gain} before encountering a distinct diminishing-returns elbow. Our findings overturn the universality of decode dominance and provide an empirical foundation for energy-optimal multimodal deployment.
\end{abstract}
